\section{System Architecture: the LLM-vs-Rule-Engine Paradigm}
\label{sec:system-architecture}

The central design decision in S.U.R.E. is not any individual model choice. It is where final
authority over a safety-relevant decision is allowed to live. This section states that decision
precisely, formalizes the two properties it depends on, and then walks the other three subsystems
-- edge vision, RAG, and MLOps drift gating -- that inherit the same discipline at a different
layer.

\subsection{The dual-layer decision system}
\label{sec:dual-layer}

Every decision cycle produces two independent status proposals. The LLM, AQUA-1B, receives the
fused sensor-and-vision state and returns a status label (\texttt{ok}/\texttt{warning}/
\texttt{critical}) plus a free-text \texttt{reasoning} field explaining it. Simultaneously, the
rule engine (\texttt{backend/rules.py}, defined in Section~\ref{sec:background}) computes its own
status label directly from the same sensor and vision thresholds, by a call path that shares no
runtime state with the LLM at all. A single function, \texttt{apply\_rule\_override}, then compares
the two: if the LLM's proposed severity is lower than the rule engine's, the final severity is
raised to the rule engine's value and the rule engine's reasoning is appended; the LLM's status is
never allowed to override the rule engine downward, and the rule engine's own decision path is
never blocked or delayed by the LLM being slow, malformed, or unavailable.

\begin{formalization}[Severity monotonicity (escalate-only invariant)]
\label{form:monotonicity}
% SOURCE_CLAIM: research_proposal.md line 61; empirically grounded by EXP02
% (experiment_workspace/experiment_runs/EXP02/results.md, module-identity check, 8/8 scenario
% agreement) and EXP03 (experiment_workspace/experiment_runs/EXP03/results.md, behavioral trace:
% no scenario's final severity fell below the rule-only computation).
Let $\preceq$ be the total order $\texttt{ok} \prec \texttt{warning} \prec \texttt{critical}$ on the
system's finite severity lattice. Writing $\mathrm{sev}_\mathrm{rule}$ for the rule engine's
computed severity and $\mathrm{sev}_\mathrm{llm}$ for the LLM's proposed severity, the system's
final decision satisfies
\[
\mathrm{sev}_\mathrm{final} \;=\; \max_{\preceq}\bigl(\mathrm{sev}_\mathrm{rule},\,
\mathrm{sev}_\mathrm{llm}\bigr).
\]
This is not a proved theorem; it is a restatement, in one line, of what
\texttt{apply\_rule\_override} already does in code, offered as a formal name for the
``escalate-only, never downgrade'' invariant the rest of the paper refers back to. Its correctness
as a \emph{safety} property depends entirely on one further condition that is not visible in the
formula itself: that \texttt{backend/rules.py} is the single source computing
$\mathrm{sev}_\mathrm{rule}$ everywhere it is consumed, with no independently maintained copy to
drift out of sync. That is exactly the property this system's own history violated once -- an
evaluation harness once carried its own copy of the rule logic, and the copy disagreed with
the system's real decision path on the single scenario that matters most
(\texttt{fish\_count==0}: the copy said \texttt{warning}, the real decision path said \texttt{ok})
-- and exactly the property
Section~\ref{sec:results} verifies directly, at the level of Python object identity, not
output-matching.
\end{formalization}

This distinction between \emph{how} a threshold is escalated and \emph{whether the free-text
justification for it can be trusted} is the boundary the second formalization exists to name.

\begin{formalization}[Enumerable vs.\ non-enumerable failure modes]
\label{form:enumerable}
% SOURCE_CLAIM: research_proposal.md line 63; empirically grounded by EXP03
% (experiment_workspace/experiment_runs/EXP03/results.md: 6/8 scenarios contain reasoning that
% materially misrepresents the true sensor snapshot, invisible to the status-only override; 3 of
% those 6 (T02/T05/T08) fabricate a specific numeric sensor value never present in the input,
% including 2 cases inside the "agrees" bucket -- see Section~5.1).
An \emph{enumerable} failure mode is a finite, coded, checkable predicate over sensor or state
variables -- for instance, dissolved oxygen falling below 6~mg/L. A \emph{non-enumerable} failure
mode is an open-ended natural-language judgment: a causal narrative that may be internally
incoherent or factually wrong about the input it claims to describe, while still landing on the
categorically correct status label. Formalization~\ref{form:monotonicity}'s escalate-only override
is a sound, verifiable safety net for the former, by construction -- it only ever reads the
enumerable \texttt{status} field. It provides \emph{no} protection whatsoever against the latter,
because \texttt{apply\_rule\_override} never inspects the \texttt{reasoning} field at all. This
distinction is not a hedge added after the fact; it bounds exactly what
Formalization~\ref{form:monotonicity} is entitled to claim, and Section~\ref{sec:results} reports
the first concrete, in-system evidence that both failure classes actually occur on this system's
own scenario set, not merely that they are conceivable.
\end{formalization}

\paragraph{AQUA-1B and adapter provenance.} AQUA-1B is a Gemma-3-1B-class model fine-tuned with a
LoRA domain adapter for this system's Turkish-language reasoning task. It is worth stating plainly,
at first mention, what that adapter actually is: the adapter loaded in the running configuration
(\texttt{sure-aqua-adapter}) was trained on eight handwritten examples, and the project's own
fine-tuning script documents this exact dataset as insufficient -- \emph{``8 handwritten examples,
smoke-test only''} -- while a properly sized, 128-example alternative dataset sat unused on disk,
predating this adapter's own training run. This is disclosed here, not as an apology, but as the honest scope of
what is measured: every AQUA-1B result reported for the decision layer in this paper
(Section~\ref{sec:results}) reflects this smoke-test adapter, not the untested larger one. Framed
as a stress test rather than a limitation, this is the more informative condition to have measured
first -- if the architecture's safety property holds against a model given close to no fair chance
to reason well, that is stronger evidence for the property's independence from model quality than
a result obtained against a carefully tuned model would have been. This scope note applies strictly
to results derived from this adapter (Section~\ref{sec:results}'s decision-layer findings); it does
not apply to the agentic tool-routing results reported in the same section, which use the base
model with no adapter loaded at all.

\subsection{Edge vision pipeline}
\label{sec:vision-pipeline}

Detection runs YOLOv11s over each incoming frame, with ByteTrack \citep{zhang2022bytetrack}
maintaining per-fish tracks across frames so that a single occluded or re-detected fish is not
double-counted. The detector was trained on 510 labeled images (412 train / 98 validation, single
class \texttt{sturgeon}). These are real photographs, captured from the project's own physical
camera and system setup and hand-labeled by the team -- not synthetic or scraped imagery -- but the
dataset is intentionally small at this stage, and expanding it is planned as future work, not
treated here as already sufficient. The detector runs via CoreML on the Apple Neural Engine for
on-device inference; Section~\ref{sec:experimental-setup} details the training and export methodology, and
Section~\ref{sec:results} reports the accuracy and latency measurements across every export format
this system was benchmarked on. The configured confidence threshold that gates which detections are
kept is a separate design choice from the threshold that maximizes F1 on a held-out curve, and the
two do not coincide here -- a fact this paper's own re-measurement surfaced rather than assumed.

\begin{formalization}[Vision operating-point tradeoff framing]
\label{form:operatingpoint}
% SOURCE_CLAIM: synthesized in formalized_results.md's G7 section from
% experiment_workspace/experiment_runs/EXP07/results.md (0.782 deployed vs. 0.719
% headline-argmax); not present verbatim in research_proposal.md prior to EXP07 -- elevated here
% because it recurs as an emergent finding, not merely a planned deliverable.
A detector's configured confidence threshold selects one specific point on its precision--recall
curve. That point is, in general, distinct from the threshold that maximizes F1 on a held-out
curve -- the two coincide only by construction or by coincidence, never by default. For a
downstream binary safety rule that consumes recall directly (here, \texttt{fish\_count==0}
triggering a \texttt{warning}), the quantity that governs operational risk is the recall
\emph{at the configured threshold}, not the recall at the academic-style F1-argmax that a validation
report typically headlines. Section~\ref{sec:results} shows this is not a merely theoretical
distinction for this system: the two numbers differ materially (0.782 versus 0.719), and citing
the wrong one understates, not overstates, the system's safety margin.
\end{formalization}

\subsection{Retrieval-augmented generation}
\label{sec:rag-pipeline}

The RAG layer indexes eight domain-specific Turkish-language documents in \texttt{pgvector},
chunked by document heading rather than by a fixed word count, and embedded with \texttt{e5-small}
\citep{wang2022e5}. A query is answered by retrieving the highest-similarity chunks above a
minimum-similarity threshold and supplying them to the LLM as grounding context with an inline
citation marker; below the threshold, nothing is retrieved and the model reasons without external
context. The threshold is set to 0.85 in the running configuration, one step above the empirical F1-optimum of
0.84 (Section~\ref{sec:results}), a deliberate choice to favor precision -- suppressing a marginal,
possibly-irrelevant retrieval entirely -- over recall, because a fabricated, confidently cited
context is judged more costly in this domain than a missed one the rule engine can still act on
independently. Tool routing for the agent layer described in Section~\ref{sec:results} is
deterministic code, not LLM reasoning, for the same underlying reason: correctness of \emph{which}
tool is called is treated as an enumerable decision the system should not delegate to a model whose
tool-selection reliability has not been established.

\subsection{MLOps and drift-gated retraining}
\label{sec:mlops}

Because frames processed at inference time carry no ground-truth labels, drift in the vision model
is inferred from the model's own output: the distribution of detection confidence scores, compared
against a reference distribution captured when the reference window was first established, via the
Population Stability Index
(Section~\ref{sec:background}). A drift signal alone only says the world changed; it says nothing
about whether a retrained candidate would actually be better, so a second, independent gate stands
between a drift alert and a registry update.

\begin{formalization}[PSI three-way gate]
\label{form:psigate}
% SOURCE_CLAIM: research_proposal.md line 64; empirically grounded by EXP08
% (experiment_workspace/experiment_runs/EXP08/results.md: no-drift window -> none, PSI=0.0048;
% significant-drift window -> retrain, PSI=0.9070; gate()'s three cases all behave as designed).
Given a reference distribution, quantile-edge (decile) binning of that reference, and a PSI value
computed over a sliding current window against those bins, the system maps the resulting PSI value
to a three-way decision $\{\texttt{none}, \texttt{review}, \texttt{retrain}\}$ at fixed thresholds
(\texttt{MODERATE}=0.10, \texttt{SIGNIFICANT}=0.25). A \texttt{retrain} decision does not by itself
promote a new model: any retrained candidate must additionally beat the incumbent by more than a
fixed \texttt{MIN\_IMPROVEMENT} margin (0.005 mAP50, absolute) before it is registered, a guard
against shipping a delta indistinguishable from ordinary training noise. Section~\ref{sec:results}
reports both halves of this gate exercised directly -- the drift decision and the promotion gate --
and reports a genuine correction to how sensitively the two candidate binning schemes (quantile
versus equal-width) actually behave across a severity sweep, rather than assuming either
convention's textbook behavior.
\end{formalization}

The choice between quantile and equal-width binning is not a cosmetic implementation detail: the
detector's confidence scores cluster narrowly (roughly 0.5--0.9), so a small number of equal-width
bins can each hold a large, uneven share of the reference mass, and Section~\ref{sec:results}
traces exactly how that structural property changes PSI's sensitivity as a function of drift
severity.
